\documentclass[11pt]{article}

\usepackage[T1]{fontenc}
\usepackage[utf8]{inputenc}
\usepackage[margin=1in]{geometry}
\usepackage{microtype}
\usepackage{amsmath,amssymb,amsthm}
\usepackage{graphicx}
\usepackage{booktabs}
\usepackage{xcolor}
\usepackage{listings}
\usepackage[hidelinks]{hyperref}
\usepackage[nameinlink,noabbrev]{cleveref}

\newtheorem{theorem}{Theorem}
\newtheorem{definition}{Definition}

\title{A Clear and Specific Paper Title}
\author{First Author \\ Institution \\ \texttt{author@example.edu}}
\date{}

\begin{document}

\maketitle

\begin{abstract}
State the problem, method, principal result, and implication in a compact
paragraph.
\end{abstract}

\section{Introduction}
\label{sec:introduction}

Explain the research question and the evidence that motivates it. Cite related
work with ordinary BibTeX keys, for example \cite{lamport1994latex}.

\section{Method}
\label{sec:method}

Define the method precisely. A numbered equation can be referenced without
hard-coding its number:
\begin{equation}
  \label{eq:objective}
  \mathcal{L}(\theta) = \frac{1}{n}\sum_{i=1}^{n}
  \ell\bigl(f_{\theta}(x_i), y_i\bigr).
\end{equation}

\section{Results}
\label{sec:results}

Report uncertainty and experimental conditions alongside point estimates.

\begin{table}[t]
  \centering
  \caption{Replace this table with the primary comparison.}
  \label{tab:results}
  \begin{tabular}{lrr}
    \toprule
    Method & Accuracy & Time (s) \\
    \midrule
    Baseline & 0.81 & 14.2 \\
    Proposed & 0.86 & 12.8 \\
    \bottomrule
  \end{tabular}
\end{table}

As \cref{eq:objective,tab:results} illustrate, labels keep references stable
when the paper changes.

\section{Conclusion}

Summarize the supported claim, limitations, and next experiment.

\bibliographystyle{plain}
\bibliography{references}

\end{document}
